% Part: model-theory
% Chapter: models-of-arithmetic
% Section: introduction

\documentclass[../../../include/open-logic-section]{subfiles}

\begin{document}

\olfileid{mod}{mar}{int}
\section{పరిచయం}

అంకగణితపు \emph{ప్రామాణిక నమూనా} అంటే
$\Domain{N} = \Nat$ గల \tetoken{నిర్మాణం}{!!{structure}}~$\Struct{N}$;
దీనిలో $\Obj{0}$, $\prime$, $+$, $\times$, $<$ లకు మనం
నిరీక్షించే అర్థనిర్దేశాలే ఉంటాయి. అంటే, $\Obj{0}$ అనేది $0$,
$\prime$ అనేది ఉత్తరాధికారి ప్రమేయం, $+$ అనేది సంకలనం,
$\times$ అనేది $\Nat$లోని సంఖ్యల గుణకారం. మరింత నిర్దిష్టంగా,
\begin{align*}
  \Assign{\Obj{0}}{N} & = 0\\
  \Assign{\prime}{N}(n) & = n + 1\\
  \Assign{+}{N}(n, m) & = n + m\\
  \Assign{\times}{N}(n, m) & = nm
\end{align*}
అయితే $\Lang{L_A}$కు $\Nat$ కాక వేరే
\tetoken{వ్యక్తి క్షేత్రాలు}{!!{domain}s} గల నిర్మాణాలు కూడా ఉన్నాయి.
ఉదాహరణకు, $\Domain{M}=\{a\}^*$ (ఒక్క $a$ సంకేతంతో ఏర్పడే పరిమిత
అనుక్రమాలు, అంటే $\emptyset$, $a$, $aa$, $aaa$, \dots) గల
$\Struct{M}$ను తీసుకొని, అర్థనిర్దేశాలను ఇలా ఇవ్వవచ్చు:
\begin{align*}
  \Assign{\Obj{0}}{M} & = \emptyset\\
  \Assign{\prime}{M}(s) & = s \concat a\\
  \Assign{+}{M}(a^n, a^m) & = a^{n + m}\\
  \Assign{\times}{M}(a^n, a^m) & = a^{nm}
\end{align*}
\[
\Assign{<}{M} =
  \Setabs{\tuple{a^n,a^m}}{n,m\in\Nat \text{ మరియు } n<m}
\]
\sourcecorrection{OLTEMODARI-001}{$\{a\}^*$లోని వాదనలు
$a^n,a^m$ కావలసి ఉండగా ఆంగ్ల మూల సూత్రాలు వాటిని $n,m$గా
రాస్తాయి; ప్రమేయాల రకాలు సరిపడేలా వాదనలను సరిచేశాం.}
\sourcecorrection{OLTEMODARI-002}{$\Struct{M}$ను
$\Lang{L_A}$-నిర్మాణంగా పేర్కొన్నప్పటికీ ఆంగ్ల మూలం $<$కు
అర్థనిర్దేశం ఇవ్వదు; $\Struct{N}$తో సమరూపత నిజమయ్యేలా పదాల పొడవుల
క్రమాన్ని స్పష్టంగా చేర్చాం.}
ఈ రెండు నిర్మాణాలు ``మౌలికంగా ఒకటే'':
వాటి \tetoken{వ్యక్తి క్షేత్రాల}{!!{domain}s}
\tetoken{మూలకాలు}{!!{element}s} మాత్రమే వేరు; అర్థనిర్దేశ ప్రమేయాల
ద్వారా \tetoken{ఆ మూలకాలు}{!!{element}s} పరస్పరం సంబంధించబడే తీరు
వేరు కాదు. ఈ రెండు
\tetoken{నిర్మాణాలు}{!!{structure}s} \emph{సమరూపాలు} అని అంటాం.

సంహతత్వ సిద్ధాంతానికి సులభ పర్యవసానంగా, $\Struct{N}$లో సత్యమైన
ఏ సిద్ధాంతానికైనా $\Struct{N}$తో సమరూపం కాని నమూనాలు కూడా ఉంటాయి.
అలాంటి నిర్మాణాలను \emph{అప్రామాణిక} నిర్మాణాలు అంటారు. వాటిలోని
ఆసక్తికరమైన విషయం ఇది: ఒక ప్రామాణిక నమూనాలోని
\tetoken{మూలకాలు}{!!{element}s} (అంటే $\Struct{N}$లోనివి మాత్రమే
కాక, దానితో సమరూపమైన అన్ని \tetoken{నిర్మాణాల్లోని}{!!{structure}s}
మూలకాలు కూడా) ప్రామాణిక సంఖ్యాపదాలు $\num{n}$ తీసుకునే విలువలతో
పూర్తిగా అయిపోతాయి, అంటే
\[
\Domain{N} = \Setabs{\Value{\num{n}}{N}}{n \in \Nat}
\]
అప్రామాణిక నమూనాల్లో అలా కాదు: $\Struct{M}$ అప్రామాణికమైతే, ప్రతి
$n$కూ $x\neq\Value{\num{n}}{M}$ అయ్యే కనీసం ఒక
$x\in\Domain{M}$ ఉంటుంది.

ఈ అప్రామాణిక మూలకాలు ఆసక్తికరమైనవి: అవి ``అనంత సహజ సంఖ్యలు.''
అదే సమయంలో వాటి ఉనికి అసంపూర్ణత పరిణామాలను ఒక అర్థంలో వివరిస్తుంది.
ఉదాహరణకు పియానో అంకగణితపు సంగతత్వ వాక్యం $\OCon[\Th{PA}]$ను,
అంటే $\lnot\lexists[x][\OPrf[\Th{PA}](x,\gn{\lfalse})]$ను
పరిశీలించండి. $\Th{PA}$ అటు $\OCon[\Th{PA}]$ను గాని, ఇటు
$\lnot\OCon[\Th{PA}]$ను గాని నిరూపించదు; కాబట్టి వాటిలో దేనినైనా
$\Th{PA}$కు సంగతంగా చేర్చవచ్చు. $\Th{PA}$ సంగతమైనది కాబట్టి
$\Sat{N}{\OCon[\Th{PA}]}$, అందువల్ల
$\Sat/{N}{\lnot\OCon[\Th{PA}]}$. కనుక $\Struct{N}$ అనేది
$\Th{PA}\cup\{\lnot\OCon[\Th{PA}]\}$కు నమూనా కాదు; దాని నమూనాలన్నీ
అప్రామాణికాలే కావాలి. $\Th{PA}\cup\{\lnot\OCon[\Th{PA}]\}$ నమూనాలో
$\lexists[x][\OPrf[\Th{PA}](x,\gn{\lfalse})]$ను సత్యం చేసే
సాక్షిగా పనిచేసే ఒక \tetoken{మూలకం}{!!{element}} ఉండాలి; అంటే అది
$\Th{PA}$ నుంచి వైరుధ్యానికి గల
\tetoken{ఒక వ్యుత్పత్తి}{!!a{derivation}} యొక్క గోడెల్ సంఖ్య కావాలి.
\sourcecorrection{OLTEMODARI-003}{సాక్షి $x$ను ప్రవేశపెట్టిన వెంటనే
వచ్చే ఆంగ్ల మూలంలోని రెండవ నిరూపణ-విధేయ ప్రయోగం ఆ వాదనను వదిలేస్తుంది;
ముందు ఇచ్చిన సంగతత్వ సూత్రానికీ వాక్యార్థానికీ సరిపడేలా
$x$ను పునఃస్థాపించాం.}
అలాంటి \tetoken{మూలకం}{!!a{element}} ప్రామాణికం కాలేదు---ఎందుకంటే
ప్రతి $n$కూ
$\Th{PA}\Proves\lnot\OPrf[\Th{PA}](\num{n},\gn{\lfalse})$.

\end{document}
